\section{Diseño de la Metaheurística para la Resolución del Problema}

Dado que el problema de distribución de combustibles combina decisiones de ruteo, asignación de productos a compartimentos, capacidades, estabilidad de carga y ventanas de tiempo, su resolución exacta llega a ser compleja al aumentar el tamaño de la instancia. Por esta razón, se propone utilizar una estrategía que consiste en usar una metaheurística para llegar a una resolución aproximada, manteniendo como base la formulación matemática determinista presentada anteriormente.


\subsection{Metaheurística Seleccionada}

La metaheurística seleccionada corresponde a \textit{Simulated Annealing} o Recocido Simulado. Esta técnica se basa en un proceso iterativo de búsqueda local, en el cual se parte desde una solución inicial y a lo largo de las iteraciones se generan soluciones vecinas mediante pequeñas modificaciones. 

Si una solución vecina mejora el costo de la solución actual, esta se acepta directamente. En caso contrario, puede ser aceptada con una cierta probabilidad que depende de la temperatura actual del algoritmo. Este mecanismo permite evitar que el algoritmo se encierre en óptimos locales durante las primeras etapas de la búsqueda y concentrarse progresivamente en soluciones de mejor calidad a medida que la temperatura disminuye.

\subsection{Representación de una Solución}

Cada solución candidata representa una planificación completa de distribución. Para ello, se define una solución $S$ compuesta por los siguientes elementos:

\begin{itemize}
    \item La asignación de estaciones de servicio a cada camión.
    \item El orden de visita de las estaciones dentro de la ruta de cada camión.
    \item La asignación de productos a los compartimentos de cada camión.
    \item Las cantidades de combustible Regular y Diésel entregadas en cada estación.
    \item Las horas de llegada estimadas a cada estación.
    \item La carga restante en cada compartimento después de cada entrega.
\end{itemize}

De forma general, una solución puede representarse como:

\begin{equation*}
S = \{R_k, A_k, Q_k\}_{k \in K}
\end{equation*}

Donde:

\begin{itemize}
    \item $R_k$: Ruta asignada al camión $k$.
    \item $A_k$: Asignación de productos a los compartimentos del camión $k$.
    \item $Q_k$: Cantidades entregadas por el camión $k$ en las estaciones de su ruta.
\end{itemize}

Por ejemplo, una solución posible puede estar dada por:

\begin{equation*}
R_{T1} = [1,3], \qquad R_{T2} = [2,4]
\end{equation*}

Lo anterior indica que el camión $T1$ visita las estaciones 1 y 3, mientras que el camión $T2$ visita las estaciones 2 y 4.

\subsection{Implementación de la Metaheurística}

Para aplicar la metaheurística de \textit{Simulated Annealing} al problema de distribución de combustibles, fue necesario transformar cada elemento definido en el modelo matemático en estructuras de datos que permitieran representar soluciones, generar movimientos vecinos y evaluar su calidad de manera eficiente.

La implementación considera una representación explícita de los camiones, estaciones de servicio, compartimentos y demandas, permitiendo construir soluciones factibles, calcular tanto las penalizaciones asociadas al combustible no entregado y asegurar que la solución cumpla con las restricciones detalladas. 

\subsubsection{Representación de los Camiones}

Cada camión se modela como una estructura que contiene la información necesaria para realizar las entregas de combustible. Para cada vehículo se almacenan:

\begin{itemize}
\item Identificador del camión.
\item Capacidad de cada compartimento.
\item Producto asignado a cada compartimento.
\item Ruta asociada al vehículo.
\item Carga disponible en cada compartimento.
\end{itemize}

De esta manera, la capacidad total del vehículo queda determinada por la suma de las capacidades de sus compartimentos, mientras que las restricciones de compatibilidad se controlan mediante la asignación de un único producto por compartimento.

\subsubsection{Representación de las Estaciones de Servicio}

Cada estación se representa mediante un conjunto de atributos que describen sus requerimientos operacionales:

\begin{itemize}
\item Identificador de la estación.
\item Coordenadas o matriz de distancias.
\item Demanda de combustible Regular.
\item Demanda de combustible Diésel.
\item Inicio de la ventana de tiempo.
\item Fin de la ventana de tiempo.
\end{itemize}

Esta información permite calcular las cantidades entregadas, los tiempos de viaje y las posibles penalizaciones por incumplimiento de horarios.

\subsubsection{Representación de una Solución}

Una solución candidata se implementa como una colección de rutas asociadas a los camiones disponibles. Cada ruta contiene la secuencia de estaciones que serán visitadas y la información de las entregas realizadas.

Conceptualmente, una solución puede representarse como:

\begin{equation}
S = \{R_k, A_k, Q_k\}_{k \in K}
\end{equation}

donde:

\begin{itemize}
\item $R_k$ corresponde a la secuencia de estaciones visitadas por el camión $k$.
\item $A_k$ representa la asignación de combustibles a los compartimentos del camión.
\item $Q_k$ almacena las cantidades entregadas en cada estación.
\end{itemize}

Además, durante la evaluación se calculan de forma dinámica los tiempos de llegada, la carga remanente en cada compartimento y los indicadores utilizados para verificar la estabilidad de carga.

\subsubsection{Proceso General del Algoritmo}

La ejecución comienza generando una solución inicial. Posteriormente, en cada iteración del Recocido Simulado, se selecciona aleatoriamente un operador de vecindad para generar una nueva solución.

La solución vecina es evaluada considerando costos de transporte, utilización de camiones, demanda no satisfecha y penalizaciones por incumplimiento de restricciones. Dependiendo de la diferencia de costo y de la temperatura actual, la nueva solución puede ser aceptada o rechazada.

Este procedimiento se repite hasta alcanzar el criterio de término definido, obteniéndose finalmente la mejor solución encontrada durante el proceso de búsqueda.

\subsection{Función de Evaluación}

La calidad de una solución se mide mediante una función de evaluación que considera los costos definidos en el modelo matemático y penalizaciones por incumplimiento de restricciones. De esta forma, el algoritmo puede determinar la factibilidad de las soluciones generadas, favoreciendo aquellas que tengan menor costo y menor violación de restricciones.

La función de evaluación se define como:

\begin{equation}
f(S) = C_{dist}(S) + C_{fijo}(S) + C_{shortage}(S) + P_{inf}(S)
\end{equation}

Donde:

\begin{itemize}
    \item $C_{dist}(S)$: Costo total asociado a la distancia recorrida.
    \item $C_{fijo}(S)$: Costo fijo de los camiones utilizados.
    \item $C_{shortage}(S)$: Penalización por demanda no satisfecha.
    \item $P_{inf}(S)$: Penalización por incumplimiento de restricciones.
\end{itemize}

El costo por distancia se calcula como:

\begin{equation}
C_{dist}(S) = c_d \sum_{k \in K} Distancia(R_k)
\end{equation}

El costo fijo se calcula según los camiones utilizados:

\begin{equation}
C_{fijo}(S) = \sum_{k \in K} F_k y_k
\end{equation}

La penalización por demanda no satisfecha se define como:

\begin{equation}
C_{shortage}(S) = c_s \sum_{j \in N} \sum_{p \in P} s_{jp}
\end{equation}

Finalmente, la penalización por infactibilidad agrupa las violaciones a las restricciones principales del problema:

\begin{equation}
P_{inf}(S) = \lambda_1 P_{cap}(S) + \lambda_2 P_{est}(S) + \lambda_3 P_{time}(S) + \lambda_4 P_{comp}(S)
\end{equation}

Donde $\lambda_1, \lambda_2, \lambda_3$ y $\lambda_4$ son factores de penalización suficientemente grandes.

\subsection{Penalización por Restricciones}

La metaheurística evalúa cada solución verificando el cumplimiento de las principales restricciones del problema. En caso de incumplimiento, se agrega una penalización a la función de evaluación.

\begin{itemize}
    \item \textbf{Penalización por capacidad:} Se aplica cuando la cantidad transportada en un compartimento supera su capacidad máxima.
    \begin{equation}
    P_{cap}(S) = \sum_{k \in K} \sum_{c \in C} 
    \max \left\{0, carga_{kc} - Q_{kc} \right\}
    \end{equation}

    \item \textbf{Penalización por estabilidad de carga:} Se aplica cuando la diferencia entre los \textit{fill ratios} de los compartimentos de un mismo camión supera el valor permitido $\Delta = 0.30$.
    \begin{equation}
    P_{est}(S) = \sum_{k \in K} \sum_{j \in R_k}
    \max \left\{0,
    \left| \frac{L_{kC_0j}}{Q_{kC_0}} -
    \frac{L_{kC_1j}}{Q_{kC_1}} \right| - \Delta
    \right\}
    \end{equation}

    \item \textbf{Penalización por ventanas de tiempo:} Se aplica cuando un camión llega antes del inicio o después del término de la ventana de tiempo de una estación.
    \begin{equation}
    P_{time}(S) = \sum_{k \in K} \sum_{j \in R_k}
    \left[
    \max \{0, a_j - t_{jk} \}
    +
    \max \{0, t_{jk} - b_j \}
    \right]
    \end{equation}

    \item \textbf{Penalización por compatibilidad:} Se aplica cuando un compartimento intenta transportar o entregar más de un tipo de combustible durante la misma ruta.
    \begin{equation}
    P_{comp}(S) =
    \sum_{k \in K} \sum_{c \in C}
    \max \left\{0, nprod_{kc} - 1 \right\}
    \end{equation}
\end{itemize}

\subsection{Generación de la Solución Inicial}

La solución inicial, osea, con la que se empieza a construir mejores soluciones, se construye asignando las estaciones disponibles a los camiones de manera simple y luego definiendo un orden inicial de visita. Para esto, se puede utilizar una estrategia constructiva basada en cercanía al depósito o una asignación inicial definida manualmente.

En una primera versión, se puede considerar la solución propuesta por el operador como solución inicial:

\begin{equation*}
R_{T1} = [1,3], \qquad R_{T2} = [2,4]
\end{equation*}

Luego, la metaheurística intenta mejorar esta solución mediante movimientos vecinos que modifican la asignación de estaciones, el orden de visita o la asignación de productos a compartimentos.

\subsection{Movimientos Vecinos}

Para explorar el espacio de soluciones, se definen distintos movimientos que generan una nueva solución $S'$ a partir de una solución actual $S$.

\begin{itemize}
    \item \textbf{Intercambio de estaciones entre camiones:} Se selecciona una estación de la ruta de un camión y se intercambia con una estación de otro camión.
    
    \item \textbf{Movimiento de una estación:} Se mueve una estación desde la ruta de un camión hacia la ruta de otro camión.
    
    \item \textbf{Cambio de orden de visita:} Se modifica el orden de las estaciones dentro de la ruta de un mismo camión.
    
    \item \textbf{Cambio de asignación de productos:} Se intercambia la asignación de productos entre los compartimentos de un camión, por ejemplo, asignar Regular a $C_1$ y Diésel a $C_0$.
\end{itemize}

Estos movimientos permiten explorar diferentes combinaciones de rutas y asignaciones de carga, manteniendo relación directa con las variables de decisión del modelo matemático.

\subsection{Criterio de Aceptación}

Sea $S$ la solución actual y $S'$ una solución vecina. Se calcula la diferencia de costo:

\begin{equation}
\Delta f = f(S') - f(S)
\end{equation}

Si la solución vecina mejora el costo, es decir, si $\Delta f < 0$, se acepta directamente:

\begin{equation*}
S \leftarrow S'
\end{equation*}

Si la solución vecina no mejora el costo, puede ser aceptada con probabilidad:

\begin{equation}
P(aceptar) = e^{-\frac{\Delta f}{T}}
\end{equation}

Donde $T$ representa la temperatura actual del algoritmo. A temperaturas altas, existe mayor probabilidad de aceptar soluciones peores, lo que permite explorar más ampliamente el espacio de búsqueda. A medida que la temperatura disminuye, el algoritmo se vuelve más selectivo, esto con el fin de no quedarse en optimos locales en las iteraciones iniciales y poder encontrar la mejor solución de un espacio de busqueda más grande.

\subsection{Esquema de Enfriamiento}

La temperatura se actualiza mediante un esquema geométrico:

\begin{equation}
T \leftarrow \alpha T
\end{equation}

Donde $\alpha$ es un parámetro de enfriamiento tal que:

\begin{equation*}
0 < \alpha < 1
\end{equation*}

En este trabajo, se pueden utilizar parámetros iniciales como:

\begin{itemize}
    \item Temperatura inicial: $T_0 = 1000$.
    \item Temperatura final: $T_f = 0.01$.
    \item Factor de enfriamiento: $\alpha = 0.95$.
    \item Número de iteraciones por temperatura: $100$.
\end{itemize}

Estos valores pueden ajustarse experimentalmente según el comportamiento observado durante las ejecuciones.

\subsection{Criterio de Término}

El algoritmo finaliza cuando se cumple alguna de las siguientes condiciones:

\begin{itemize}
    \item La temperatura alcanza un valor menor o igual a la temperatura final $T_f$.
    \item Se alcanza un número máximo de iteraciones.
    \item No se observa mejora en la mejor solución durante un número determinado de iteraciones consecutivas.
\end{itemize}

Al finalizar, el algoritmo retorna la mejor solución encontrada durante todo el proceso de búsqueda.

\subsection{Pseudocódigo de la Metaheurística}

\begin{lstlisting}
Generar solucion inicial S
Evaluar f(S)
S_mejor = S
T = T_inicial

Mientras T > T_final:
    Para iteracion = 1 hasta iteraciones_por_temperatura:
        
        S_vecina = generar_vecino(S)
        Evaluar f(S_vecina)
        
        delta = f(S_vecina) - f(S)
        
        Si delta < 0:
            S = S_vecina
        Sino:
            r = numero aleatorio entre 0 y 1
            Si r < exp(-delta / T):
                S = S_vecina
        
        Si f(S) < f(S_mejor):
            S_mejor = S
    
    T = alpha * T

Retornar S_mejor
\end{lstlisting}

\subsection{Validación de la Solución Obtenida}

Una vez obtenida la mejor solución, se realiza una validación final sin considerar penalizaciones, verificando directamente el cumplimiento de las restricciones del problema. Esta validación considera:

\begin{itemize}
    \item Que las rutas comiencen y terminen en el depósito.
    \item Que cada estación sea atendida como máximo por un camión.
    \item Que la demanda sea satisfecha o que el shortage sea correctamente calculado.
    \item Que cada compartimento transporte un único producto.
    \item Que no se excedan las capacidades de los compartimentos.
    \item Que la diferencia entre fill ratios no supere $\Delta = 0.30$.
    \item Que las llegadas a las estaciones respeten sus ventanas de tiempo.
\end{itemize}

De esta manera, la metaheurística permite obtener una solución candidata y luego comprobar formalmente si dicha solución cumple las condiciones operacionales del problema.

